

\begin{figure}[H]
    \centering
    \fbox{%
        \parbox{0.95\textwidth}{
            %\raggedright % Makes the content left-aligned
            %\raggedright % Makes the content left-aligned
            \ttfamily % Monospaced font
            %\fontfamily{cmtt}\selectfont % Use a monospaced font
            \fontsize{8pt}{10pt}\selectfont % Set font size
            %\setlength{\baselineskip}{8pt} % Adjust line spacing here
            Your task is to extract data on wild bee observations described in the report passages into structured JSONs. The expected json format is specified below. Extract only the specified attributes. If an attribute of an observation is not specified or you can't find information for this attribute, leave it out. Don't make up information! Stick strictly to the information in the report passage. Extract information exactly as it appears in the passage.\\

            ATTRIBUTE DESCRIPTIONS:\\
            bee: Latin designation of the bee. Consisting of genus and species\\
            taxon: name of the author who first described the bee, and sometimes the respective publication year. In most cases the taxon is located right behind the latin designation of the bee\\
            n\_males: number of observed male bees, could be abbreviated with "\male"-symbol after the respective amount\\
            n\_females: number of female bees/ worker bees, could be abbreviated with "\female"-symbol after the respective amount\\
            n\_queens: the number of queens is explicitly described\\
            n\_divers: number of bees with unassignable gender, could be abbreviated with "\mercury"-symbol after the respective amount\\
            leg: person who confirmed the bee, could be abbreviated with "leg." in front of the respective name\\
            coll: person who collected the bee, could be abbreviated with "coll." in front of the respective name\\
            det: person who has determined the bee, could be abbreviated with "det." or "vid." in front of the respective name\\

            All attributes in their defined order:\\
            bee, taxon, date, location, n\_males, n\_females, n\_queens, n\_divers, leg, coll, det, habitat, visited\_flowers, behaviour, observed\_nesting, collecting\_method, remarks\\
            
            EXPECTED JSON FORMAT:\\
            
            The JSON shall contain a list of observations as value for the key "observations". Each observation should be a dictionary containing the sorted attributes specified above as keys and observation-specific values extracted from the report passage. If an attribute is not specified for an observation, leave it out. If an attribute has the same value for all observations, it should be reported as meta-attribute with the prefix "meta\_". Meta-attributes should be specified before "observations" in the order defined above. If an attribute is specified as a meta-attribute, it should not be repeated in the observation-dictionary. Only "location" can be specified as meta-attribute and observation-specific attribute at the same time. If you can't find any wild bee observations in the report passage return an empty JSON.}
        }%
    \caption{Extraction-Prompt consisting of initial task instruction, explanation of attributes, order of further attributes and description of the output-format}
\end{figure}